\documentclass[11pt]{article}

% Change "review" to "final" to generate the final (sometimes called camera-ready) version.
% Change to "preprint" to generate a non-anonymous version with page numbers.
\usepackage[final]{acl}

\usepackage{times}
\usepackage{latexsym}

\usepackage[T1]{fontenc}

\usepackage[utf8]{inputenc}

\usepackage{microtype}

\usepackage{inconsolata}

\usepackage{graphicx}

\usepackage{booktabs}

\usepackage{rotating}

\usepackage{url}

\title{NLP Assignment 2}

\author{ Author \\
  MTU / (John) Paul Nagle \\
  \texttt{John.P.Nagle@mymtu.ie} \\}


\begin{document}
\maketitle
\begin{abstract}
This paper presents a cascaded speech-to-text translation system for the low-resource Irish-to-English language pair, developed as part of the Natural Language Processing 
module for the M.Sc. in AI at MTU. It addresses the challenges of translating Irish (Gaeilge) speech while implementing a two-stage pipeline combining Automatic Speech 
Recognition (ASR) and Machine Translation (MT) components. The baseline system integrates OpenAI's Whisper model \citep{radford2023whisper} for ASR with Meta's NLLB-200 
\citep{nllb2022} for neural machine translation. To address the low-resource nature of Irish, we develop an improved system utilising Irish-specific fine-tuned Whisper 
models that better handle the language's unique phonological characteristics. We evaluate both systems on the IWSLT 2026 Irish-English test set using standard metrics 
including BLEU \citep{papineni2002bleu} and chrF++ \citep{popovic2015chrf}, with comprehensive error analysis identifying common failure modes such as language misidentification 
and morphological errors. Our results demonstrate the effectiveness of language-specific model adaptation in low-resource settings, while highlighting persistent challenges 
in cascaded architectures including error propagation from ASR to MT stages.
\end{abstract}

\section{Problem Statement}

This work addresses the challenge of low-resource speech-to-text translation for the Irish-to-English language pair. Speech translation (ST) — converting spoken audio in one 
language into text in another — presents significant difficulties for low-resource languages like Irish (Gaeilge), which has limited transcribed audio data and fewer high-quality 
pretrained models compared to well-resourced languages.

We implement a cascaded pipeline architecture combining Automatic Speech Recognition (ASR) and Machine Translation (MT) components. This approach offers practical advantages 
in low-resource settings: ASR and MT components can be independently pretrained on large datasets, errors can be diagnosed at each stage, and strong off-the-shelf models are 
readily available. Our baseline system integrates OpenAI's Whisper model for ASR with Meta's NLLB-200 for neural machine translation.

The primary research question is: \textit{How can we improve translation quality for low-resource Irish speech by optimizing component selection in a cascaded pipeline?} 
We evaluate two key dimensions: (1) the impact of Irish-specific fine-tuned ASR models versus general multilingual models, and (2) the comparative performance of multilingual 
(NLLB-200) versus bilingual (OPUS-MT) translation models. Our evaluation uses the IWSLT 2026 Irish-English dataset \citep{iwslt2026dataset}, measuring performance through 
BLEU, chrF++, and COMET metrics while conducting systematic error analysis to identify failure modes specific to Irish morphology and phonology.

\section{Dataset}

We utilize the IWSLT 2026 Irish-English speech translation dataset \citep{iwslt2026dataset}, which provides parallel Irish speech and English text for low-resource speech 
translation research. For this work, we exclusively use the development split located in \texttt{iwslt2026\_ga-eng/iwslt2025\_ga-eng/dev/}, as it is the only subset providing 
both Irish audio files (\texttt{wav/}) and corresponding English reference translations (\texttt{txt/}).

\subsection{Dataset Composition}

The development set contains 1120 audio samples with corresponding English translations. This limited dataset size reflects the low-resource nature of Irish speech translation 
and necessitates reliance on pretrained models rather than training from scratch. While the dataset includes additional splits (training and test sets), these were not 
utilised in our pipelines due to the absence of paired English translations required for evaluation.

\subsection{Audio Characteristics}

All audio files are provided in WAV format with 16kHz sampling rate, mono channel configuration. Audio segments vary in duration, with filenames set so that an alphabetical listing
matches the order of the translations in the dev.eng. The dataset includes both individual utterances and slightly longer continuous speech segments, reflecting realistic single sentance 
speech translation scenarios, but not longer conversations.

\subsection{Annotation Format}

Ground truth annotations are provided in TSV format (\texttt{stamped.tsv}) containing three fields: audio filename, start timestamp, and length in seconds. 
In some of the other unused data folders, english reference texts are also available as plain text files (\texttt{train.eng}, \texttt{dev.eng}, \texttt{test.eng}), but in the dev 
set we only have the dev.eng file for translations.

\subsection{Preprocessing}

Audio preprocessing is minimal to preserve authentic speech characteristics. We load audio files using the \texttt{librosa} library with 16kHz sampling rate to match 
Whisper model requirements. No additional noise reduction, normalization, or voice activity detection is applied, as pretrained models are designed to handle raw speech input. 
For translation evaluation, we tokenize reference texts using standard whitespace splitting to compute BLEU and chrF++ metrics.

\subsection{Dataset Challenges}

The dataset presents several challenges characteristic of low-resource speech translation: (1) extremely limited training data (1120 samples) necessitating reliance on 
pretrained models, (2) Irish phonological complexity including lenition, eclipsis, and dialectal variation, and (3) morphological richness with extensive inflection and mutation 
systems.


\section{Pipeline Architectures}

Baseline and improved pipeline architectures with justifications.

\section{Experiment}

\subsection{Evaluation Setup}

Footnotes are inserted with the \verb|\footnote| command.\footnote{This is a footnote.}

\subsection{Metrics}

metrics

\subsection{Experimental Conditions}

\subsection{Citations}

You can use the command \verb|\citep| (cite in parentheses) to get ``(author, year)'' citations \citep{Gusfield:97}.

\citep{Gusfield:97}

\section{BLEU/chrF++ Scores}

\begin{table*}[!]
\centering
\caption{Overall Comparison Table}
\label{tab:system-comparison}
\footnotesize
\begin{tabular}{@{}lrrrrrrrrrrr@{}}
\toprule
\textbf{System} & \rotatebox{90}{\textbf{Samples}} & \rotatebox{90}{\textbf{Valid}} & \rotatebox{90}{\textbf{BLEU}} & \rotatebox{90}{\textbf{chrF++}} & \rotatebox{90}{\textbf{COMET}} & \rotatebox{90}{\textbf{Coverage}} & \rotatebox{90}{\textbf{Rep \%}} & \rotatebox{90}{\textbf{Avg Hyp Len}} & \rotatebox{90}{\textbf{Avg Ref Len}} & \rotatebox{90}{\textbf{Length Ratio}} & \rotatebox{90}{\textbf{Avg Time/File (s)}} \\
\midrule
Specialised ASR+NLLB & 1120 & 1120 & 4.36 & 22.25 & 0.4752 & 100.0 & 2.59 & 10.40 & 9.88 & 1.05 & 0.82  \\
Specialised ASR+OPUS-MT & 1120 & 1120 & 3.40 & 21.56 & 0.4384 & 100.0 & 0.62 & 10.55 & 9.88 & 1.07 & 0.64  \\
Baseline ASR+NLLB & 1120 & 1120 & 1.51 & 15.90 & 0.4375 & 100.0 & 4.11 & 10.31 & 9.88 & 1.04 & 1.70  \\
Baseline ASR+OPUS-MT & 1120 & 1120 & 1.42 & 15.19 & 0.4168 & 100.0 & 2.32 & 10.07 & 9.88 & 1.02 & 1.56  \\
\bottomrule
\end{tabular}
\end{table*}

\begin{table*}[!]
\centering
\caption{Relative Comparison vs Baseline ASR + NLLB}
\label{tab:relative-comparison}
\footnotesize
\begin{tabular}{@{}lrrrrrrrr@{}}
\toprule
\textbf{System} & \rotatebox{90}{\textbf{BLEU}} & \rotatebox{90}{\textbf{BLEU $\Delta$ vs Baseline}} & \rotatebox{90}{\textbf{chrF++}} & \rotatebox{90}{\textbf{chrF++ $\Delta$ vs Baseline}} & \rotatebox{90}{\textbf{Coverage}} & \rotatebox{90}{\textbf{Coverage $\Delta$ vs Baseline}} & \rotatebox{90}{\textbf{Avg Time/File (s)}} & \rotatebox{90}{\textbf{Time/File $\Delta$ vs Baseline}} \\
\midrule
Specialised ASR + NLLB & 4.36 & 2.85 & 22.25 & 6.35 & 100.0 & 0.0 & 0.82 & -0.88 \\
Specialised ASR + OPUS-MT & 3.40 & 1.89 & 21.56 & 5.66 & 100.0 & 0.0 & 0.64 & -1.06 \\
Baseline ASR + NLLB & 1.51 & 0.00 & 15.90 & 0.00 & 100.0 & 0.0 & 1.70 & 0.00 \\
Baseline ASR + OPUS-MT & 1.42 & -0.09 & 15.19 & -0.71 & 100.0 & 0.0 & 1.56 & -0.14 \\
\bottomrule
\end{tabular}
\end{table*}




\section{Qualitative Analysis}

\section{Result Interpretations}

\section{Summary}

% Bibliography entries for the entire Anthology, followed by custom entries
%\bibliography{anthology,custom}
% Custom bibliography entries only
\bibliography{custom}

\appendix

\section{AI Use Declaration}
\label{sec:appendix}

I used AI to blah blah blah.

\end{document}